\begin{abstract}
Models that predict deterioration in the emergency department (ED) are often evaluated on the whole stay, although part of that record is produced after the decision the score should inform, and their uncertainty is seldom examined for patients who deteriorate. In this study, we audited an end-of-stay deterioration model on MIMIC-IV-ED, a retrospective single-center cohort, for input leakage, and redefined the task at a landmark two hours after arrival, with intensive care admission or in-hospital death within the next 24 hours as the outcome. On a chronological, patient-level split, we compared a gated recurrent unit with decay (GRU-D), gradient-boosted trees, and logistic regression with equal tuning, and evaluated static, adaptive, and class-conditional conformal prediction sets on the \RVHtwoTestEligible{} test-period stays eligible at the landmark. On the end-of-stay task, arrival information alone reached an area under the receiver operating characteristic curve (AUROC) of \RVIncArr{}, length of stay and measurement counts added \RVIncWorkOverArr{}, and the full model reached \RVEosHybridAuroc{}. At the landmark, timing-aware trees reached \RVAurocXgbT{} and exceeded the GRU-D model by \RVTreesOverHybrid{} (95\% CI \RVTreesOverHybridLo{}--\RVTreesOverHybridHi{}). Prediction sets with near-nominal marginal coverage covered only \RVSetCovOneStaticLac{} of the patients with the outcome. Class-conditional sets raised this to \RVSetCovOneMondEci{}, but the set containing only the deterioration label had a positive predictive value of \RVSetPpvMondEci{}, with \RVSetFaMondEci{} false alarms per 100 patients, matching a plain threshold on the score. Deterioration models should be evaluated at a stated prediction time, and conformal sets should be reported with class-specific coverage and false-alarm burden.
\end{abstract}
